\section{n2. Full invariant-boundary fibres are group orbits}
\textit{Supports: ENTIRE qualified G331; n1.}

Let $S$ be an excellent reduced local ring of positive dimension, with a finite group $G$ acting
by local automorphisms, $R=S^G$ excellent local, and $S$ finite over $R$.
The domain and formal instances of these hypotheses are established in n3.
No characteristic, perfection, faithful-action or invertibility of $|G|$
is needed in this theorem. For every primary $\mathfrak a\subset R$, let
$Y=V_{\mathfrak aS}(S)$, $X=V_{\mathfrak a}(R)$ and $\rho:Y\to X$.
By G331, this is a compact surjective restriction map. We prove its
EVERY fibre, including ALL nonminimal supports, is exactly a $G$ orbit.

For a primary/unit $S$ ideal $J$, choose finite $R$-module generators
$b_1,\ldots,b_l$ of $J$. Form the orbit polynomial
\[
 N_J(T)=\prod_{\sigma\in G}\left(\sum_i T_i\sigma(b_i)\right)\in R[T].
\]
Its coefficients belong to $R$ because $G$ permutes the factors.
Let $\mathcal N_G(J)$ be their ideal in $R$. A change of module generators
is a linear substitution over $R$; substitutions in both directions show
this coefficient ideal is independent of the generators. Redundant
module generators may therefore be added freely. If $J\supset\mathfrak
 a^N S$ and $x$ is a generator of $\mathfrak a^N$, add $x$ among the
$b_i$; its pure coefficient is $x^{|G|}$. These finitely many powers have
radical $\mathfrak m_R$, so the coefficient ideal is primary or unit.
It is proper for proper $J$: at a normalized lift all the finitely many
$(w\circ\sigma)(J)$ are strictly positive, so the coefficient value below
is positive, whereas a unit coefficient ideal would have value zero.
For $J=S$ adding $1$ proves $\mathcal N_G(J)=R$.

For every lift $w$ with arbitrary prime support $P$, first pass to the
valued field $\operatorname{Frac}(S/P)$; coefficients killed by $P$ are
zero there. In this field the minimum coefficient value of a nonzero
polynomial is multiplicative. Divide each factor by a coefficient of
least value and reduce modulo the valuation maximal ideal: the two
nonzero residue polynomials have nonzero product. This proves the Gauss
identity without an assumption on the support or on a residue extension.
Each factor here has a finite least value because $J$ contains a primary
power. Thus
\[
 v(\mathcal N_G(J))=\sum_{\sigma\in G}(w\circ\sigma)(J),\qquad v=\rho w.
\]
In particular the finite orbit measure $\sum_\sigma\delta_{w\circ\sigma}$
has integrals on ALL primary/unit gauges determined by $v$ alone.

For completeness these gauges determine finite measures on $Y$.
Their rational difference span is a lattice after common denominators,
using
\[
 \min(w(I)-w(J),w(K)-w(H))=w(IH+KJ)-w(JH).
\]
It contains rational constants from $w(\mathfrak aS)=1$, and separates
EVERY pair of points using $w((b)+(\mathfrak aS)^N)=\min(w(b),N)$.
Its uniform closure is a real vector lattice. Here is its compact
approximation proof, so no new Stone--Weierstrass black box is needed.
For a continuous function $f$ and each $x,y$, an affine transform of a
separating function matches $f$ at $x,y$ (a constant works if $x=y$).
Fix $x$; finitely many $y$-neighbourhoods and the maximum of their matching
functions give a function everywhere above $f-\epsilon$ and equal to
$f$ at $x$. It is below $f+\epsilon$ near $x$. Finitely many such
$x$-neighbourhoods and their minimum give uniform error at most
$\epsilon$. Thus the closure is all $C^0(Y)$.

Any two lifts of $v$ therefore have equal orbit measures: the preceding
identity gives equality on the difference lattice and uniform density
extends equality to every continuous function, since both measures have
mass $|G|$. Their finite supports must coincide. Indeed, if a point of
one support were absent from the other, compact Hausdorff normality gives
a continuous function equal to one there and zero on the other finite
support, contradicting equality and positive integer mass. Each lift
occurs with positive multiplicity in its own orbit measure. Consequently
\[
 \rho^{-1}(v)=G\cdot w\quad\text{for EVERY }v\in X.
\]
No generic degree calculation or valuation-extension transitivity theorem
has been imported. Existence of lifts is G331's quotient-support lifting:
kill the support first, choose a prime above it, extend the valuation of
that quotient fraction field through integral closure and a dominating
valuation ring, and use the algebraic value-group real embedding.
The argument here then proves transitivity on ALL such choices, including
support-prime and valuation extensions simultaneously. It never extends
a nonzero-support semivaluation to the original fraction field of $S$.

Since $\rho$ is compact surjective, the induced continuous bijection
$Y/G\to X$ is a homeomorphism. The quotient is open, because the inverse
image of the image of an open set $U$ is $\bigcup_{\sigma\in G}\sigma U$.
This topological statement alone does not supply positivity descent;
that descent is the ENTIRE G331 theorem. Its consequence is the exact
identification of $C_R,D_R$ with the invariant parts of $C_S,D_S$.
The orbit coefficient construction also explicitly maps a rational
negative primary profile to its orbit average:
\[
 -v(\mathcal N_G(J))/(|G|n)
   =|G|^{-1}\sum_\sigma[-(w\circ\sigma)(J)/n].
\]
The left side has an actual negative effective ample presentation on its
own primary blowup. Units give zero. This is a boundary-cone construction,
not a Monge--Ampere product or a claim about unrestricted nonnormal traces.

Artifacts: notes/n2.tex and P12.3-quotients-and-wild-conductors.tex.
The finite diagnostic polynomial is $(Tx+Uy)(-Tx+Uy)=U^2y^2-T^2x^2$;
it checks a small exact instance only. Rerun: python exact\_checks.py;
python build\_delivery.py; python verify\_delivery.py.
